% Written last, after every section is approved. Every figure in it must trace to the numeric evidence ledger.
\noindent
US structured note issuers print an estimated value on the cover of every
pricing supplement, computed with models and a funding rate the filing does not
disclose. This paper asks whether an outside party can reproduce that figure,
and tests the question on the hardest instrument it covers: worst-of
autocallable notes on three correlated equity indices, whose redemption before
maturity is the issuer's to initiate rather than automatic on five of the six
notes studied. Surfaces
are fitted to public option quotes under no-arbitrage constraints, converted to
local volatility, and simulated jointly, with the issuer's call solved by
least-squares Monte Carlo from the issuer's side of the contract.

The disclosed figure is not reproduced. Independent valuations land below it on
all six notes, by \$11.41 to \$27.92 per \$1{,}000 of principal across the five
whose barriers are covered by quoted strikes, against simulation standard
errors below \$0.40. Solving for the funding spread that would reconcile the
two valuations returns a negative spread on every note, from $-278$ to $-123$
basis points against an A-rated corporate benchmark of 69, so the undisclosed
funding rate cannot account for the difference. Four further candidates are
excluded on direction or magnitude. Two survive, the correlation assumption and
the shape of the volatility surface, and a third, the dynamics of that surface,
lies outside what a local-volatility model can test; none of the three is
published in the filings. A
census of the same week bounds the exercise: of 1{,}086 supplements filed, 129
are final notes of this type, and 71 of those can be priced from public quotes
at all.
